\section{Literature Positioning and Gap}

Weak-form efficiency and return-predictability studies show that apparent forecast gains are often unstable once out-of-sample discipline is imposed \parencite{lomackinlay1988,welchgoyal2008,campbellthompson2008,rapach2010}. In emerging markets, integration and information-friction mechanisms motivate context-specific reassessment rather than universal claims \parencite{harvey1995,bekaertharvey1995}. For Latin America, much of the evidence remains concentrated in variance-ratio, adaptive-market, or related dependence-testing traditions \parencite{urrutia1995,sanchezgranero2020,lo2004,cruzhernandez2024}, while large-scale machine-learning evidence is mostly developed-market focused \parencite{gu2020}. Diffusion methodology is technically mature in machine learning \parencite{ho2020,song2021}, and finance applications are emerging \parencite{cho2026}, but benchmark-disciplined weak-form tests in this LatAm ETF setting remain limited.

The gap is therefore not the absence of EMH tests, but the absence of a coherent benchmark-first forecast-comparison design that separates predictive association from structural inefficiency claims. This paper addresses that gap by combining leakage-audited target construction, expanding-window evaluation, strong baseline comparators, formal loss-differential inference, and explicit interpretation boundaries for EMH.
